\documentclass{article}
\usepackage[utf8]{inputenc}
\usepackage[citestyle=authortitle-terse,backref]{biblatex}
\addbibresource{bibliography.bib}


\usepackage[left=4cm, right=4cm]{geometry}
\usepackage{palatino,eulervm,dutchcal,xcolor}
\usepackage{graphicx,subcaption,float}
\usepackage{enumitem,parskip,multicol,array}
\newcolumntype{L}{>{$}c<{$}}%table in mathmode
\usepackage{amsthm,amssymb,amsmath,mathtools,thmtools,tikz,tikz-cd}
\usetikzlibrary{%
	matrix,%
	calc,%
	arrows,%
	shapes,
	decorations.markings,backgrounds,calc,intersections}
\tikzcdset{scale cd/.style={every label/.append style={scale=#1},
		cells={nodes={scale=#1}}}}
\usepackage[bookmarks,bookmarksopen,bookmarksdepth=3]{hyperref}
\hypersetup{%colores
	colorlinks=true,
	urlcolor=blue,
	linkcolor=magenta,
	citecolor=blue,
	filecolor=blue,
	urlbordercolor=white,
	linkbordercolor=white,
	citebordercolor=white,
	filebordercolor=white}
\usepackage{cleveref}

\definecolor{blue-violet}{rgb}{0.54, 0.17, 0.89}
\definecolor{azure}{rgb}{0.0, 0.5, 1.0}
\definecolor{green(ncs)}{rgb}{0.0, 0.62, 0.42}
\definecolor{forestgreen}{rgb}{0.13, 0.55, 0.13}
\definecolor{limegreen}{rgb}{0.2, 0.8, 0.2}
\definecolor{palatinateblue}{rgb}{0.15, 0.23, 0.89}
\definecolor{trueblue}{rgb}{0.0, 0.45, 0.81}
\definecolor{goldenyellow}{rgb}{1.0, 0.87, 0.0}
\definecolor{fashionfuchsia}{rgb}{0.96, 0.0, 0.63}
\definecolor{brightcerulean}{rgb}{0.11, 0.67, 0.84}
\definecolor{jonquil}{rgb}{0.98, 0.85, 0.37}
\definecolor{lavendermagenta}{rgb}{0.93, 0.51, 0.93}
\definecolor{peru}{rgb}{0.8, 0.52, 0.25}
\definecolor{persimmon}{rgb}{0.93, 0.35, 0.0}
\definecolor{persianred}{rgb}{0.8, 0.2, 0.2}
\definecolor{persianblue}{rgb}{0.11, 0.22, 0.73}
\definecolor{persiangreen}{rgb}{0.0, 0.65, 0.58}
\definecolor{persianyellow}{rgb}{0.9, 0.89, 0.0}



\declaretheoremstyle[headfont=\color{trueblue}\normalfont\bfseries,]{colored1}
\declaretheoremstyle[headfont=\color{forestgreen}\normalfont\bfseries,]{colored2}
\declaretheoremstyle[headfont=\color{peru}\normalfont\bfseries,]{colored3}
\declaretheoremstyle[headfont=\color{persiangreen}\normalfont\bfseries,]{colored4}
\declaretheoremstyle[headfont=\color{brightcerulean}\normalfont\bfseries,]{colored5}
\declaretheoremstyle[headfont=\color{lavendermagenta}\normalfont\bfseries,]{colored6}
\declaretheoremstyle[headfont=\color{blue-violet}\normalfont\bfseries,]{colored7}
\declaretheoremstyle[headfont=\color{green(ncs)}\normalfont\bfseries,]{colored8}
\declaretheoremstyle[headfont=\color{peru}\normalfont\bfseries,]{colored9}
\declaretheoremstyle[headfont=\color{persiangreen}\normalfont\bfseries,]{colored10}

\declaretheorem[style=colored1,numbered=no,name=Theorem]{thm}
\declaretheorem[style=colored2,numbered=no,name=Proposition]{prop}
\declaretheorem[style=colored3,numbered=no,name=Lemma]{lemma}
\declaretheorem[style=colored4,numbered=no,name=Corollary]{coro}
\declaretheorem[style=colored5,numbered=no,name=Example]{example}
\declaretheorem[style=colored5,numbered=no,name=Examples]{examples}
\declaretheorem[style=colored7,numbered=no,name=Exercise]{exercise}
\declaretheorem[style=colored7,numbered=no,name=Conjecture]{conjecture}
\declaretheorem[style=colored6,numbered=no,name=Remark]{remark}
\declaretheorem[style=colored8,numbered=no,name=Claim]{claim}
\declaretheorem[style=colored9,numbered=no,name=Definition]{defn}
\declaretheorem[style=colored10,numbered=no,name=Question]{question}

\numberwithin{equation}{section}

\newcommand{\A}{\mathbb{A}}
\newcommand{\R}{\mathbb{R}}
\newcommand{\Z}{\mathbb{Z}}
\newcommand{\N}{\mathbb{N}}
\newcommand{\C}{\mathbb{C}}
\newcommand{\Q}{\mathbb{Q}}
\newcommand{\D}{\mathbb{D}}
\renewcommand{\P}{\mathbb{P}}

\newcommand{\Ac}{\mathcal{A}}
\newcommand{\Bc}{\mathcal{B}}
\newcommand{\Cc}{\mathcal{C}}
\newcommand{\Dc}{\mathcal{D}}
\newcommand{\Ec}{\mathcal{E}}
\newcommand{\Fc}{\mathcal{F}}
\newcommand{\Gc}{\mathcal{G}}
\newcommand{\Hc}{\mathcal{H}}
\newcommand{\Ic}{\mathcal{I}}
\newcommand{\Jc}{\mathcal{J}}
\newcommand{\Kc}{\mathcal{K}}
\newcommand{\Lc}{\mathcal{L}}
\newcommand{\Mc}{\mathcal{M}}
\newcommand{\Nc}{\mathcal{N}}
\newcommand{\Oc}{\mathcal{O}}
\newcommand{\Pc}{\mathcal{P}}
\newcommand{\Qc}{\mathcal{Q}}
\newcommand{\Rc}{\mathcal{R}}
\newcommand{\Sc}{\mathcal{S}}
\newcommand{\Tc}{\mathcal{T}}
\newcommand{\Uc}{\mathcal{U}}
\newcommand{\Vc}{\mathcal{V}}
\newcommand{\Wc}{\mathcal{W}}
\newcommand{\Xc}{\mathcal{X}}
\newcommand{\Yc}{\mathcal{Y}}
\newcommand{\Zc}{\mathcal{Z}}

\newcommand{\mf}{\mathfrak{m}}

\DeclareMathOperator{\img}{img}
\DeclareMathOperator{\Arg}{Arg}
\DeclareMathOperator{\id}{id}
\DeclareMathOperator{\pt}{pt}
\DeclareMathOperator{\Alt}{Alt}
\DeclareMathOperator{\sgn}{sgn}
\DeclareMathOperator{\hTop}{h-Top}
\DeclareMathOperator{\supp}{supp}
\DeclareMathOperator{\Int}{Int}
\DeclareMathOperator{\sing}{sing}
\DeclareMathOperator{\Ob}{Ob}
\DeclareMathOperator{\Mor}{Mor}
\DeclareMathOperator{\Top}{Top}
\DeclareMathOperator{\Set}{Set}
\DeclareMathOperator{\CGWH}{CGWH}
\DeclareMathOperator{\Hom}{Hom}
\DeclareMathOperator{\Map}{Map}
\DeclareMathOperator{\Tot}{Tot}
\DeclareMathOperator{\op}{op}
\DeclareMathOperator{\ev}{ev}
\DeclareMathOperator{\hofib}{hofib}
\DeclareMathOperator{\rel}{rel}
\DeclareMathOperator{\Nat}{Nat}
\DeclareMathOperator{\Arr}{Arr}
\DeclareMathOperator{\Funct}{Funct}
\DeclareMathOperator{\colim}{colim}
\DeclareMathOperator{\odd}{odd}
\DeclareMathOperator{\GL}{GL}
\DeclareMathOperator{\SL}{SL}
\DeclareMathOperator{\SO}{SO}
\DeclareMathOperator{\U}{U}
\DeclareMathOperator{\SU}{SU}
\DeclareMathOperator{\Sp}{Sp}
\DeclareMathOperator{\M}{M}
\DeclareMathOperator{\Aut}{Aut}
\DeclareMathOperator{\PGL}{PGL}
\DeclareMathOperator{\PSL}{PSL}
\DeclareMathOperator{\Vect}{Vect}
\DeclareMathOperator{\VectBund}{VectBund}
\DeclareMathOperator{\Open}{Open}
\DeclareMathOperator{\coker}{coker}
\DeclareMathOperator{\con}{con}
\DeclareMathOperator{\Hol}{Hol}
\DeclareMathOperator{\Fl}{Fl}
\DeclareMathOperator{\Gr}{Gr}
\DeclareMathOperator{\Td}{Td}
\DeclareMathOperator{\Bl}{Bl}
\DeclareMathOperator{\Pic}{Pic}

\begin{document}
	{\Large\bfseries Fano manifolds and Mirror symmetry }
	
	
	\subsection*{Fano manifolds}
	\begin{itemize}
		\item Algebraic geometry
		
		Let $X$ be a compact closed complex submanifold of $\C\P^N$.
		\[X\hookrightarrow \C\P^N\]
		Also, $H=\C\P^{N-1}$ is a submanifold of $\C\P^N$ so consider the intersection $X\cap H$. Also we have an ample bundle,
		\[c_1(i^*(\Oc(1)))=H^1\sim c_1(X)=c_1(Tx)\]
		
		\item Differential geometry
		
		Also we have a symplectic form $\omega$ such that
		\begin{enumerate}
			\item $[\omega]=c_1(X)\in H^2(X,\Z)$.
			\item $\omega$-K\"hler, so $d\omega=0$ and $\omega(u,Ju)>0$ for all $u\neq 0$ and $\omega>0$.
		\end{enumerate}
		Also we have that $\Hol_g\subset \U(N)$.
		
		\item Symplectic geometry
		
		We have that $\omega$ is compatible with the complex structure $J$ and also we have a Riemannian metric. Let's define a manifold $M$ to be \textbf{\textit{monotone}} if $c_1(M,\omega)=\lambda[\omega]$ for some $\lambda\in\R$ positive. This is related to the ample line bundle.
		
		This is a \textit{rigid object}. (So we expect to find something rigid on the other side of the mirror.)
	\end{itemize}
	\begin{examples}\leavevmode
		\begin{enumerate}
			\item A point.
			\item Product of Fano manifolds.
			\item $\P^1\subset\A$ given by $x=y^{-1}$.
			\item $\frac{dx}{x}=-\frac{dy}{y}\implies \omega^*_x=\Oc(+2)$.
			\item $\P^n$.
			\item Hypersurface of degree $d$, $X_d\subset\P^n$. Solutions of polynomials.
			\begin{itemize}
				\item Degree 1 is just hypersurfaces.
				\item Degree 2, quadric.
				\item Degree 3, cubic. Famous work of1972 Griffiths and Clamens. They invented the technique of intermediate jacobians, …
				\item Degree 4. They are birational. Iskovskisk-Manin.
			\end{itemize}
			\item Grassmanians $\Gr(k,N)$. Is an example of $G/P$, P for parabolic and it contains Borel.
			\item $\Fl(1,2,\ldots,N)$, so $W_1\subset W_2\subset\ldots W_N$, which is an example of $G/B$, B for Borel.
			\item Fano manifolds are very important in Minimal Model Problem (MMP).
		\end{enumerate}
	\end{examples}
	
	Properties of a Fano manifold $F$:
	\begin{enumerate}
		\item $\pi_1(F)=0$.
		\item $H^i(F,\Oc_F)=H^k(F,\omega\otimes\omega^\wedge=0$. This implies that Todd genus $\chi(F,\Oc)=1=\sum(-1)\dim H^k(F,\Oc)=\int_{[x]}\Td X$ and the covering $\tilde{D}\overset{d}{\to} F$ is also a fano manifold.
		\item All this means that hodge numbers $h^{k,0}=h^{0,k}=0$.
		\item $\Pic F\cong \Z^\rho$ where $\rho=b_2$, so it is torsion free.
		\[\begin{tikzcd}
			0\arrow[r]&\Pic F\arrow[r,"c_1"]&H^2(F,\Z)\arrow[r]&0
		\end{tikzcd}\]
		\item Now consider the open complex cone $\Kc$ of a K\"hler manifold, it is a complicated object. For Fano, this cone is rational polyhedral---this follows from famous cone theorem: there are two duals, second cohomology, $H^2(X,\R)=\R^\rho$, and second homology, $H_2(X,\R)$, so if you have a complex curve $C\subset X$ its class is $[C]\in N_1$. Look for curves that $c_1>0$, so $(-\Kc x)\cdot c>0$ so on the side of cohomology you have the K\"hler cone of $\Kc$ of ample classes  and on the side homology there's the Mory cone $\Mc=\langle[c]\rangle$ and it turns out that the two cones are dual. So one is polyhedral if the other is.
		
		[…] then for every ray there is a map $F\to F_\rho$ such that it contracts all the rays in this class and only them, universal property.
		
		If you have any space $Y$, there exists so-called MRC fibration so that spaces are \textbf{\textit{rationally connected}}: for every two points you can draw rational curves. For example, for the projective space we get a point on the base.
		
		\item $F$ is arc-rationally connected. Of course, it is path-connected. Theorem by KKM, 1990.
		
		\item Deformation of Fano manifold is Fano due to positivity of the bundle (?).
		
		\item $\Gamma(X,\Omega^{\otimes N})=0$.
	\end{enumerate}
	
	\begin{thm}[Very important, KMM]
		For any dimension $d\geq0$ there exist only finitely many deformation classes of Fano manifolds.
	\end{thm}
	This yields:
	\begin{itemize}
		\item $d=0$, point
		\item $d=1$, $\P^1$.
		\item $d=2$, the are 10 \textbf{\textit{del Pezzo}} surfaces.
		\begin{itemize}
			\item $\P^2$.
			\item $\P^1\times \P^1$ (the only one that is not a blow up.)
			\item $\Bl_{P,Q}\P^2$.
			\item $\Bl_{P_1,P_2,P_3}\P^2$.
			Let's introduce some funny number: $\int_{[X]}c_1(X)^n=d>0$ the \textbf{\textit{degree}}, it is bounded by some number $F(d)$. What is \textbf{\textit{degree of $\P^n$}}? $(-K_{\P^n})^n=(n+1)^n$.
			
			\textbf{\textit{Noether's formula}} says that 
			\[d+(\rho+2)=c_1^2+c_2=(2\chi(\Oc))\]
			so
			\[d+\rho=10\]
			And we can bound the degree: $1\leq d\leq 9$. We have
			\[\begin{tikzcd}
				9\arrow[r,no head]&8\arrow[r,no head]&7\arrow[ld,no head]\arrow[r,no head]&6\arrow[r,no head]&5\arrow[r,no head]&4\arrow[r,no head]&3\arrow[r,no head]&2\arrow[r,no head]&1\\
				&8
			\end{tikzcd}\]
			So this reminds us of the Dynkin diagram of $E_{10}$. Look at degree 3. It's a cubic surface $X_3\subset\P^3$. It has 27 lines and 27 pencils of conics and corresponds with branes and compactifications of ?. Also $E_6$ has 27 something and it's not a coincidence. In fact Cartan defined $E_6$ as the automorphisms of… using that there's 45 planes, intersecting… ?
			
			Consider double cover of $\P^2$.
			
			Also $E_7·$ is related to degree 2, they have 56 lines on them. $E_8$ to degree 1, it is very complicated of course so there's 248 lines on it.
			
			Up to degree 6 they are toric surfaces. 5 already is rigid and it has no moduli, finite group of automorphisms, symmetric group acts here.
			
			Du Val: ADE singularities.
			
			5 is $\overline{\Mc_{0,5}}$ its a moduli space. It's nice.
			
			Another property is $H^{1,1}=H^2(X,\R)$. There is a unimodular Lorenzian lattice $(H^2(X,\Z),\cap\text{ pairing})$.
		\end{itemize}
		\item $d=3$ there are 105. $b_2\to 17$ Poincaré missed it.
		\item $d=4$ nobody knows at all.
	\end{itemize}
	
	\subsection*{Mirror symetry}
	Corresponding to Fano manifolds, we have \textbf{\textit{Landau-Ginzburg models}}, which are pairs $(Y,W)$ where $Y$ is a quasi-projective variety and $W\in\Gamma(Y,\Oc_Y)$. Fantastic Laurent Polynomials.
	
	\subsection{The conifold point}
	This is the abstract of \href{http://www.mpim-bonn.mpg.de/node/5969}{a talk by Sergey in 2015 at Bonn}:
	\begin{quotation}
		Conjecture $O$ describes the geometry in the complex line of the eigenvalues $u_i$ of the operator of quantum multiplication by the first Chern class acting on the cohomology of a Fano manifold. In particular, it says that eigenvalues with maximal absolute value have multiplicity one and one of them is real and positive number $T$.
		
		Fano manifolds tend to have mirror dual Ginzburg-Landau potentials $f$, which
		tend to have a distinguished non-degenerate critical point which we name the
		\textbf{\textit{conifold point}}. Explicitly the conifold point is the unique critical point $P$ on real positive locus, and the respective critical value $T_{\con} = f(P)$ is the global minimum on the real positive locus. \textbf{In this case it is conjectured that $T_{\con}$ conicides with $T$}, that is any eigenvalue $u_i$ has absolute value at most $T_{\con}$, and that the conifold point is the unique critical point with value $T_{\con}$.
		
		In most cases the existence of the conifold point is clear and the conjectures
		can be checked to be true, however we do not know how to prove them even
		for abstract toric Fano manifolds, or complete intersections therein. These conjectures
		are basic for formulation of Gamma conjectures about the appearance of Gamma
		unction in the symplectic topology of Fano manifolds. Two references are arXiv:1404.7388
		and my joint work with Golyshev and Iritani arXiv:1404:6407.
	\end{quotation}
	Also see who quoted \textit{The conifold point}.  \href{https://scholar.google.com/scholar?start=20&hl=es&as_sdt=2005&sciodt=0,5&cites=11808564544162580683&scipsc=}{here}. And also don't forget that \href{https://en.wikipedia.org/wiki/Combinatorial_mirror_symmetry}{wiki on combinatorial mirror symmetry} exists!
	
	\subsection{the history of quantum cohomology}
	\paragraph{-1.} Before mirror symmetry, in 1985, people tried to classified vector bundles in $\C P^2$. There are \textbf{\textit{rigid bundles}}.  (See J.M. Drezet et J. Le Pother.) So in 1987 Garolov and ? (she studied at the same time as Kontsevich) called them \textbf{\textit{Markov equations}} of the form $x^2+y^2+z^2=3xys$.
	
	So, there are three independant morphisms $(\Oc(0),\Oc(1),\Oc(2))\rightsquigarrow?$, difficult to specify formally. So people introduced simmer orb decompositions, and then introduced notion of Serre functor, which is donde by consider two bundles and $R\Hom(F,S\Sigma)$ and it is just the twist canonical bundle.
	
	Then some people introduced to $A_\infty$ algebras and $L_\infty$ algebras. $A_\infty$ algebra is very important because it is decomposed as $A=A_{\ev}\oplus A_{\odd}$.

	Some other old friends like Penner, who gave a cell decomposition of the moduli space $\Mc_{g,n}\times\R^n$. And Kontsevich actually uses this on his 90's book on the Witten conjectures. Then you can get some graph cohomology out of this. Kontsevich went to Harvard in the 90's and gave two lectures.
	
	So in the 90's physical mathematics really exploded.
	
	\begin{conjecture}[Witten]
		\[\int_{[\bar\Mc_{g,n,n}]}\bigcap KdV\]
	\end{conjecture}
	
	And then also Vafa with some \textbf{\textit{quantum rings}} with multiplicities on $H^\bullet(X)$ and $X$ K\"ahler manifold taget in 2d $N=(2,2)$ QFT. \textbf{\textit{White's associative product}}. But why is it associative? Beginings of cohomology.
	
	In winter 92 Fukaya gave a lecture where he related some theories in codimensions and then a category, symplectic manifold $(Y,\omega)$. Take Lagrangian manifolds as objects and then 
	
\end{document}
